\section{Exact initial values and bounded-domain coverage}
\label{sec:enclosures}

\subsection{Rational bounds for an exponential}

The origin case is rational and should stay purely symbolic. For other rational
anchors, a small enclosure procedure is sufficient. Let $z\in\Q$ and choose
$s\ge0$ so that $r=|z|/2^s\le1/2$. For an integer $n\ge1$, set
\[
 S_n(r)=\sum_{k=0}^{n}\frac{r^k}{k!},\qquad
 T_n(r)=\frac{r^{n+1}}{(n+1)!}
             \frac{1}{1-r/(n+2)}.
\]
All denominators are strictly positive. The successive terms of the tail
starting at degree $n+1$ have ratios at most $r/(n+2)$, so
\begin{equation}
 S_n(r)\le e^r\le S_n(r)+T_n(r).
 \label{eq:exp-bounds}
\end{equation}
For $z<0$, invert this interval before the repeated squaring:
\[
 \frac{1}{S_n+T_n}\le e^{-r}\le\frac{1}{S_n}.
\]
Squaring a nonnegative interval is monotone. Repeating it $s$ times produces
bounds on $e^z$.

The implementation rounds lower endpoints down and upper endpoints up to a
fixed dyadic grid after arithmetic stages. If the denominator is $2^b$, the
rounding operations are exact integer floor and ceiling operations, not
floating-point conversions. Repeated squaring of a very small interval can
therefore produce a lower endpoint zero; this loses precision but remains
sound. It never converts an uncertain lower sign into a positive claim.

\subsection{Normalising an entire point evaluation}

Direct evaluation at a large positive anchor can create enormous intermediate
numbers even when the sign is easy. First evaluate every polynomial block
$P_\lambda(a)$ exactly, deleting zero values, and choose
\[
 M=\max_{P_\lambda(a)\ne0}(\lambda a).
\]
Then
\begin{equation}
 f(a)=e^M\sum_\lambda P_\lambda(a)e^{\lambda a-M}.
 \label{eq:normalised-evaluation}
\end{equation}
Every remaining exponential argument is nonpositive. The normalised sum has
the same sign as the original value and avoids computing an enormous common
positive scale. Its enclosure is what the receipt records.

This distinction matters when reading the JSON: the recorded lower and upper
numbers at a nonzero anchor bound a \emph{positively normalised} value, not
necessarily $f(a)$ itself. The checker recomputes that same mathematical
normalisation. An integration layer which exports those numbers as bounds on
$f(a)$ without multiplying by the common factor would be wrong.

When the normalised coefficients cancel exactly before enclosure, the result
is exact. When the enclosure straddles zero, the producer returns unknown at
its current precision. The prototype does not use a decimal tolerance, a
``nearly zero'' test, or an oracle's numerical sign as a proof.

\subsection{A reference interval cover, not a new interval tactic}

For experiments combining compact and unbounded domains, the archive includes
a simple natural-interval worker. It evaluates polynomial blocks by interval
Horner arithmetic, bounds each exponential using monotonicity and
\eqref{eq:exp-bounds}, sums the resulting intervals, and bisects when the lower
bound is negative.

Each leaf records its rational endpoints and its recomputed enclosure. An
internal node records a rational split point and two child nodes. The checker
requires that the split lie strictly inside the parent's interval and that the
children cover exactly its left and right pieces. A collection of individually
valid positive boxes is insufficient without this coverage proof.

Normalisation on an interval uses one common positive constant for the
\emph{whole leaf}. The chosen scale bounds every rate times every point in the
leaf, so the remaining exponential arguments are nonpositive. Normalising each
summand by a different factor would alter their relative magnitudes and would
be unsound. The implementation comments call out this specific failure mode.

This cover is deliberately not presented as a contribution competing with
Forge's existing polynomial Bernstein worker or with Lean interval libraries.
Its purpose is to exercise the new analytic interface, certify finite patches,
and retain exact counterpoints. Production should reuse proved interval
infrastructure rather than adopting a second unverified interval stack.

\subsection{Why subdivision cannot replace the ladder}

For $f(x)=e^x-1-x$ on $[0,h]$, natural interval evaluation contains the range
\[
 [1,e^h]-1-[0,h]=[-h,e^h-1].
\]
Every $h>0$ gives a negative lower endpoint, even though $f$ is nonnegative.
Splitting the interval leaves a first piece containing zero, and the same
obstruction persists. The problem is lost correlation between $e^x$ and $x$,
not insufficient effort on a remote part of the domain.

The differential ladder establishes the exact second-order contact at the
endpoint with rational initial values. This comparison is not a claim that
all sophisticated Taylor-model or symbolic interval methods must fail; a
method which explicitly factors the contact or uses the same derivative
information can succeed. It is a precise explanation of why this delivered
natural-interval worker returns unknown while the ladder succeeds.

Conversely, interval subdivision certifies strictly positive patches of
functions whose native ladder fails. In the experiment, the four instances
of $f_\varepsilon$ on $[0,2]$ use respectively 14, 21, 31, and 66 leaves as
$\varepsilon$ decreases through $1,1/2,1/4,1/16$. The proof obligations are
bounded-domain statements, unlike the auxiliary ladders' whole-ray statements.
Their counts must not be compared as though the workers solved the same goal.

\subsection{Negative points are useful, but have a different type}

A failed enclosure is not a counterexample. The worker records a rational
counterpoint only when a point enclosure has \emph{strictly negative upper
bound}. For $e^x-2$ on $[0,1]$, it finds such a point and the independent replay
checks its membership in the original interval and the negative upper bound.

A counterpoint receipt answers a refutation request. It does not make a
nonnegativity goal successful. Likewise a tail with negative eventual sign
provides valuable information about an impossible requested positive tail, but
its logical polarity must remain explicit. The JSON goal type and
\code{verify_goal} enforce these distinctions.
